\section{Introduction}
\label{sec:intro}
Large Language Models (LLMs) have demonstrated remarkable performance in software engineering and have been widely adopted for code generation ~\cite{aiXcoder-7B, cai2025ai, seedcoder}. Specifically, code generation aims to generate source code satisfying a given natural language requirement. The rapid advancement of model capabilities in code generation tasks has catalyzed the emergence of numerous coding agents, exemplified by industrial products such as GitHub Copilot~\cite{githubGitHubCopilot} and Cursor~\cite{cursorCursor}. These products are progressively transforming developers' workflows, making it crucial to establish a benchmark that accurately mirrors realistic development scenarios to further advance their capabilities.

In the landscape of code generation, early benchmarks primarily focused on function-level generation~\cite{humaneval}. This scope subsequently expanded to class-level code generation~\cite{classeval} and further deepened into repository-level context-aware completion tasks~\cite{li2024evocodebench}. Furthermore, to evaluate model capabilities within realistic software evolution, recent research has shifted towards feature-level development scenarios, such as FEA-Bench~\cite{li2025fea} and NoCode-bench~\cite{deng2025nocode}. Unlike completion tasks, feature-level code generation requires models to design and implement new features within complex repository contexts based on natural language requirements (e.g., functional documentation).

However, as demonstrated in \cref{tab:comparison}, existing feature-level benchmarks generally fall short of reflecting realistic software development scenarios. This discrepancy primarily stems from two critical limitations: 
\ding{182} \textbf{Unrealistic Task Inputs:} Existing benchmarks typically lower the task difficulty by enriching inputs with code hints, such as function signatures. In contrast, real-world software development initiates with only requirements. This compels models to reason through the implementation process independently, from identifying appropriate modification locations to formulating the specific functions needed to implement the feature. By providing these hints upfront, current benchmarks bypass the critical challenge of independently bridging the gap between user intent and concrete code implementation.
\ding{183} \textbf{Data Leakage (aka Data Contamination):} Current feature-level benchmarks have not been updated since their initial releases, making them static benchmarks. Since rapidly iterating models are often trained on the latest open-source repositories, these static benchmarks drawn from public repositories are susceptible to memorization. This fundamentally contradicts real-world scenarios where developers face novel, unseen problems. Consequently, such leakage compromises the evaluation's validity, making it difficult to distinguish true generalization capabilities from mere recollection of the solution.

\textbf{To align evaluation with realistic software development scenarios, we introduce FeatBench, a benchmark designed for realistic feature-level implementation scenarios.} Compared to existing feature-level code generation benchmarks, FeatBench aligns more closely with realistic development scenarios, distinguishing itself through two key contributions:
\ding{182} \textbf{Realistic Task Inputs:} Task inputs consist solely of natural language requirements (e.g., ``I want to...''). This format mirrors realistic software development by strictly isolating user intent from code hints, imposing a substantial challenge that requires agents to independently comprehend the repository context and devise a strategy to implement the feature without explicit hints.
\ding{183} \textbf{Evolving Data:} To mitigate the risk of data leakage inherent in static benchmarks, FeatBench is an evolving benchmark and will be dynamically updated every period (e.g., 6 months). Specifically, we build an automated collection pipeline, which constructs new versions of FeatBench from the latest repositories. The initial release of our benchmark comprises 157 tasks sourced from 27 actively maintained open-source GitHub repositories.

In addition to the aforementioned contributions, FeatBench ensures strict data quality guarantees and covers more domains:
\ding{184} \textbf{Rigorous Data Collection Process:} FeatBench employs a strict multi-level filtering pipeline to ensure high data quality. At the repository level, we strictly select actively maintained repositories requiring at least three formal releases and identifiable test suites to guarantee verifiability, while rigorously excluding non-production tutorials. Crucially, at the Pull Request (PR) level, we strictly confine tasks to those modifying existing functions without adding or deleting functions. This rigorous constraint ensures that developer-written test cases remain valid for evaluating agent-generated code. Finally, we validate each task to include at least one Fail-to-Pass (F2P) test case. Furthermore, our human evaluation confirms that 93.3\% of the sampled tasks possess comprehensive and unambiguous requirements.
\ding{185} \textbf{Comprehensive Test Cases:} Each task is equipped with a rigorous dual-validation suite comprising both Fail-to-Pass (F2P) and Pass-to-Pass (P2P) test cases. To guarantee evaluation robustness, each task includes an average of 1,694.6 test cases. This comprises an average of 2.2 F2P tests to verify the correctness of the new feature and 1,692.4 P2P tests derived from the repository's existing test suite. This massive scale of validation ensures that agents not only implement the required functionality but also maintain strict backward compatibility, preventing regressions in production-grade software.
\ding{186} \textbf{Diverse Application Domains:} FeatBench spans diverse domains such as AI/ML, DevOps, and Web. This diversity mirrors the vast landscape of real-world software engineering, ensuring that the benchmark accurately reflects realistic  development scenarios.

\begin{table}[t]
    \centering
    \caption{Comparison with existing code generation benchmarks. The \textbf{Task} column classifies benchmarks based on the granularity of the generated output. For instance, \textbf{Function-level Code Generation} indicates that the target output is confined to the scope of a single function, even if the input context originates from different files in a repository. FeatBench distinguishes itself by targeting \textbf{Feature Implementation} based on \textbf{Realistic Natural Language Requirements}, strictly devoid of code hints.}
    \resizebox{0.95\textwidth}{!}{
    \setlength{\tabcolsep}{10pt}
    \begin{tabular}{l|cccc}
        \toprule 
        \textbf{Benchmark} & \textbf{Date} & \textbf{Task} & \textbf{Data Curation}& \textbf{Code Hints}\\
        \midrule 
        HumanEval\cite{humaneval} & Jul, 2021 & Function-level Code Generation& Manual & Signatures  \\
        ClassEval\cite{classeval} & Aug, 2023 & Class-level Code Generation& Manual & Class Skeleton  \\
        CoderEval\cite{yu2024codereval} & Feb, 2024 & Function-level Code Generation& Manual & Signatures \\
        DevEval\cite{li2024deveval} & Aug, 2024 & Function-level Code Generation& Manual & Signatures  \\
        EvoCodeBench\cite{li2024evocodebench}& Dec, 2024 & Function-level Code Generation& Automatic & Signatures  \\
        FEA-Bench\cite{li2025fea} & Mar, 2025 & Feature-level Code Generation& Automatic& Signatures \\
        NoCode-bench\cite{deng2025nocode} & Jul, 2025 & Feature-level Code Generation& Automatic& Identifier Hints \\
        \midrule 
        \textbf{FeatBench} (Ours) & Feb, 2026 & \textbf{Feature-level Code Generation}& \textbf{Automatic} & \textbf{None}  \\
        \bottomrule
    \end{tabular}
    \label{tab:comparison} }
\end{table}

To demonstrate the utility of FeatBench, we evaluate two SOTA agent frameworks, Trae-agent and Agentless, using four SOTA LLMs, including open-source models like DeepSeek V3.1~\cite{deepseek} and proprietary models such as GPT-5~\cite{gpt-5}, Doubao-Seed-1.6~\cite{doubao}, and Qwen3-Coder-Flash~\cite{qwen3}. Following established standards~\cite{deng2025nocode, swebenchgoeslive}, we adopt the \textbf{Resolved Rate (\%)} as our primary metric. We also report the \textbf{Patch Apply Rate (\%)} and the \textbf{File-level Localization Success Rate (\%)}, alongside several auxiliary metrics. Based on our experimental results, we find that:
\ding{182} FeatBench poses a significant challenge to SOTA agents, with the top-performing configuration (Trae-agent~\cite{traeagent} + GPT-5~\cite{gpt-5}) achieving a resolved rate of only 29.94\%. We identify an apparent performance disparity between agent paradigms: autonomous, planning-based agents substantially outperform rigid, pipeline-based counterparts.
\ding{183} Agent performance is strictly constrained by repository and patch complexity. We observe a strong inverse correlation between repository complexity and performance: resolved rates reach up to 60–70\% for smaller repositories but degrade sharply to 10–30\% for those exceeding 800 files or 300k Lines of Code (LOC). Similarly, the resolved rate collapses for widespread modifications spanning more than 5 files or 50 LOC. Furthermore, consistent performance across different task creation periods confirms the absence of data leakage and validates the efficacy of our evolving benchmark strategy.
\ding{184} Regressive implementation is the predominant failure reason, accounting for 73.6\% of analyzed failure cases. This high regression rate is driven by a behavioral pattern termed aggressive implementation, where agents exhibit ``scope creep'' by proactively refactoring code or extending features beyond the explicit user intent. This behavior acts as a double-edged sword: although it often introduces defects in production-grade environments, it simultaneously allows agents to occasionally generate architectural designs that surpass human-written patches in terms of modularity and robustness.

In summary, the main contributions of this paper are as follows:
\begin{itemize}
    \item We identify critical limitations in existing feature-level code generation benchmarks, specifically their reliance on unrealistic task inputs (e.g., function signatures) and static datasets that lead to data leakage, failing to reflect realistic software development scenarios.
    \item We introduce and open-source FeatBench, a novel evolving feature-level benchmark derived from actively maintained open-source repositories. FeatBench aligns with realistic development by utilizing natural language requirements without code hints and continuously updating from latest repositories to mitigate data contamination.
    \item We conduct extensive experiments on FeatBench using SOTA agents and LLMs. 
    Our experiments reveal that current models face significant challenges in realistic feature implementation, with a maximum resolution rate of only 29.94\%. Furthermore, we uncover critical behavioral issues, including a high propensity for introducing regressions and a prevalence of ``scope creep,'' where agents diverge from user intent through over-engineering.
\end{itemize}